% GENERATED FILE. Edit canonical lessons, then run npm run generate:books.
% canonical-source-hash: fnv1a64:16451d2ef422b713
% canonical-lessons: FA-C327-emdadgar, FA-C327-urzhans, FA-C327-azhir, FA-C327-zende-mandan, FA-C327-ghariq-e-nejat

\chapter{Help arrives}
\label{ch:fa-help-arrives}
\hlchaptermodality{327}

\begin{chapteropening}
\textbf{By the end of this chapter:} \emph{I can talk about rescue workers, the emergency services, sirens, lifeguards and people who survived.}

Complete the last lesson of chapter 327: I can talk about rescue workers, the emergency services, sirens, lifeguards and people who survived.
\end{chapteropening}

\section[\texorpdfstring{\fa{امدادگر} (emdâdgar)}{امدادگر (emdâdgar)}]{\fa{امدادگر} (emdâdgar) — a rescue worker, a paramedic}

\label{lesson:FA-C327-emdadgar}



\begin{quote}
Before the new one: say the Persian for to faint, then the Persian for to come round.
\end{quote}

\subsection*{You'll want to know: \fa{امدادگر}}
\textbf{\fa{امدادگر}} (\emph{emdâdgar}) — \textquotedblleft{}a rescue worker, a paramedic\textquotedblright{}.

\begin{quote}
\fa{امدادگرها} \fa{زود} \fa{به} \fa{محل} \fa{رسیدند}.
\end{quote}

\begin{quote}
\emph{Emdâdgar-hâ zud be mahal residand.}
\end{quote}

\begin{quote}
The rescue workers soon reached the scene.
\end{quote}

\emph{emdâd} \textquotedblleft{}help, relief\textquotedblright{} + the native suffix \emph{-gar}.

\begin{cousinweb}
\emph{emdâdgar} joins \emph{emdâd} \textquotedblleft{}help, relief\textquotedblright{}, borrowed from Arabic \emph{imdād} (EtymDB lists the Arabic word), and the native Persian suffix \emph{-gar} \textquotedblleft{}one who does\textquotedblright{}.
\end{cousinweb}

\begin{grammarlens}[title={What you've built}]
One more word for telling, explaining and describing at B1.
\end{grammarlens}

\subsection*{Your turn}
\begin{itemize}
\raggedright
  \item \emph{Say it:} \emph{emdâdgar}
  \item \emph{Say it:} \emph{emdâdgar}, once more
  \item \emph{Say it:} a sentence of your own with \emph{emdâdgar}
  \item \emph{Recall:} say the Persian for to faint, then the Persian for to come round, then say \emph{emdâdgar} again
\end{itemize}

\begin{tcolorbox}[breakable,colback=teal!4,colframe=teal!35!black,title={Before you move on}]
What does \fa{امدادگر} mean? (\textquotedblleft{}A rescue worker, a paramedic\textquotedblright{}.) Say it once more.
\end{tcolorbox}
\section[\texorpdfstring{\fa{اورژانس} (urzhâns)}{اورژانس (urzhâns)}]{\fa{اورژانس} (urzhâns) — the emergency services; an emergency room}

\label{lesson:FA-C327-urzhans}



\begin{quote}
Before the new one: say the Persian for to come round, then the Persian for a rescue worker.
\end{quote}

\subsection*{You'll want to know: \fa{اورژانس}}
\textbf{\fa{اورژانس}} (\emph{urzhâns}) — \textquotedblleft{}the emergency services; an emergency room\textquotedblright{}.

\begin{quote}
\fa{مادرم} \fa{به} \fa{اورژانس} \fa{زنگ} \fa{زد}.
\end{quote}

\begin{quote}
\emph{Mâdaram be urzhâns zang zad.}
\end{quote}

\begin{quote}
My mother called the emergency services.
\end{quote}

Both the ambulance service and a hospital's emergency department. Written with \fa{ژ}.

\begin{cousinweb}
\emph{urzhâns} is borrowed from French \emph{urgence}, \textquotedblleft{}emergency\textquotedblright{} (EtymDB).
\end{cousinweb}

\begin{grammarlens}[title={What you've built}]
One more word for telling, explaining and describing at B1.
\end{grammarlens}

\subsection*{Your turn}
\begin{itemize}
\raggedright
  \item \emph{Say it:} \emph{urzhâns}
  \item \emph{Say it:} \emph{urzhâns}, once more
  \item \emph{Say it:} a sentence of your own with \emph{urzhâns}
  \item \emph{Recall:} say the Persian for to come round, then the Persian for a rescue worker, then say \emph{urzhâns} again
\end{itemize}

\begin{tcolorbox}[breakable,colback=teal!4,colframe=teal!35!black,title={Before you move on}]
What does \fa{اورژانس} mean? (\textquotedblleft{}The emergency services; an emergency room\textquotedblright{}.) Say it once more.
\end{tcolorbox}
\section[\texorpdfstring{\fa{آژیر} (âzhir)}{آژیر (âzhir)}]{\fa{آژیر} (âzhir) — a siren, an alarm}

\label{lesson:FA-C327-azhir}



\begin{quote}
Before the new one: say the Persian for a rescue worker, then the Persian for the emergency services.
\end{quote}

\subsection*{You'll want to know: \fa{آژیر}}
\textbf{\fa{آژیر}} (\emph{âzhir}) — \textquotedblleft{}a siren, an alarm\textquotedblright{}.

\begin{quote}
\fa{صدای} \fa{آژیر} \fa{آمبولانس} \fa{از} \fa{دور} \fa{میآمد}.
\end{quote}

\begin{quote}
\emph{Sedâ-ye âzhir-e âmbulâns az dur miâmad.}
\end{quote}

\begin{quote}
The sound of an ambulance siren was coming from far away.
\end{quote}

Written with \fa{ژ}. \emph{Âzhir-e khatar} is an alarm siren.

\begin{cousinweb}
The origin of \emph{âzhir} \textquotedblleft{}a siren, an alarm\textquotedblright{} is not settled, and EtymDB lists no entry.
\end{cousinweb}

\begin{grammarlens}[title={What you've built}]
One more word for telling, explaining and describing at B1.
\end{grammarlens}

\subsection*{Your turn}
\begin{itemize}
\raggedright
  \item \emph{Say it:} \emph{âzhir}
  \item \emph{Say it:} \emph{âzhir}, once more
  \item \emph{Say it:} a sentence of your own with \emph{âzhir}
  \item \emph{Recall:} say the Persian for a rescue worker, then the Persian for the emergency services, then say \emph{âzhir} again
\end{itemize}

\begin{tcolorbox}[breakable,colback=teal!4,colframe=teal!35!black,title={Before you move on}]
What does \fa{آژیر} mean? (\textquotedblleft{}A siren, an alarm\textquotedblright{}.) Say it once more.
\end{tcolorbox}
\section[\texorpdfstring{\fa{زنده} \fa{ماندن} (zende mândan)}{زنده ماندن (zende mândan)}]{\fa{زنده} \fa{ماندن} (zende mândan) — to survive, to stay alive}

\label{lesson:FA-C327-zende-mandan}



\begin{quote}
Before the new one: say the Persian for the emergency services, then the Persian for a siren.
\end{quote}

\subsection*{You'll want to know: \fa{زنده} \fa{ماندن}}
\textbf{\fa{زنده} \fa{ماندن}} (\emph{zende mândan}) — \textquotedblleft{}to survive, to stay alive\textquotedblright{}.

\begin{quote}
\fa{همه} \fa{مسافران} \fa{زنده} \fa{ماندند}.
\end{quote}

\begin{quote}
\emph{Hame-ye mosâferân zende mândand.}
\end{quote}

\begin{quote}
All the passengers survived.
\end{quote}

Literally \textquotedblleft{}to remain alive\textquotedblright{}. Present stem \emph{mân-}: \emph{zende mimânad}.

\begin{cousinweb}
\emph{zende} \textquotedblleft{}alive\textquotedblright{} is inherited from Middle Persian \emph{zīndag}, and \emph{mândan} \textquotedblleft{}to remain\textquotedblright{} is inherited from Middle Persian too (EtymDB).
\end{cousinweb}

\begin{grammarlens}[title={What you've built}]
One more word for telling, explaining and describing at B1.
\end{grammarlens}

\subsection*{Your turn}
\begin{itemize}
\raggedright
  \item \emph{Say it:} \emph{zende mândan}
  \item \emph{Say it:} \emph{zende mândan}, once more
  \item \emph{Say it:} a sentence of your own with \emph{zende mândan}
  \item \emph{Recall:} say the Persian for the emergency services, then the Persian for a siren, then say \emph{zende mândan} again
\end{itemize}

\begin{tcolorbox}[breakable,colback=teal!4,colframe=teal!35!black,title={Before you move on}]
What does \fa{زنده} \fa{ماندن} mean? (\textquotedblleft{}To survive, to stay alive\textquotedblright{}.) Say it once more.
\end{tcolorbox}
\section[\texorpdfstring{\fa{غریق} \fa{نجات} (ghariq-e nejât)}{غریق نجات (ghariq-e nejât)}]{\fa{غریق} \fa{نجات} (ghariq-e nejât) — a lifeguard}

\label{lesson:FA-C327-ghariq-e-nejat}



\begin{quote}
Before the new one: say the Persian for a siren, then the Persian for to survive.
\end{quote}

\subsection*{You'll want to know: \fa{غریق} \fa{نجات}}
\textbf{\fa{غریق} \fa{نجات}} (\emph{ghariq-e nejât}) — \textquotedblleft{}a lifeguard\textquotedblright{}.

\begin{quote}
\fa{غریق} \fa{نجات} \fa{پسر} \fa{را} \fa{از} \fa{آب} \fa{بیرون} \fa{آورد}.
\end{quote}

\begin{quote}
\emph{Ghariq-e nejât pesar râ az âb birun âvard.}
\end{quote}

\begin{quote}
The lifeguard pulled the boy out of the water.
\end{quote}

Literally \textquotedblleft{}rescue of the drowning\textquotedblright{}: \emph{ghariq} \textquotedblleft{}a drowning person\textquotedblright{} + \emph{nejât} \textquotedblleft{}rescue\textquotedblright{}.

\begin{cousinweb}
\emph{nejât} \textquotedblleft{}rescue\textquotedblright{} is borrowed from Arabic \emph{najāt} (EtymDB); \emph{ghariq} \textquotedblleft{}drowning, a drowning person\textquotedblright{} is usually traced to Arabic \emph{gharīq}. The phrase itself is Persian.
\end{cousinweb}

\begin{grammarlens}[title={What you've built}]
One more word for telling, explaining and describing at B1.
\end{grammarlens}

\subsection*{Your turn}
\begin{itemize}
\raggedright
  \item \emph{Say it:} \emph{ghariq-e nejât}
  \item \emph{Say it:} \emph{ghariq-e nejât}, once more
  \item \emph{Say it:} a sentence of your own with \emph{ghariq-e nejât}
  \item \emph{Recall:} say the Persian for a siren, then the Persian for to survive, then say \emph{ghariq-e nejât} again
\end{itemize}

\begin{tcolorbox}[breakable,colback=teal!4,colframe=teal!35!black,title={Before you move on}]
What does \fa{غریق} \fa{نجات} mean? (\textquotedblleft{}A lifeguard\textquotedblright{}.) Say it once more.
\end{tcolorbox}
